% GENERATED FILE. Edit canonical lessons, then run npm run generate:books.
% canonical-source-hash: fnv1a64:a53efaa3e3a784ed
% canonical-lessons: RU-C107-jam, RU-C107-nut, RU-C107-pepper, RU-C107-nurse, RU-C107-cook

\chapter{Jam, Nut, Pepper, Nurse, Cook}
\label{ch:ru-jam-nut-pepper-nurse-cook}
\hlchaptermodality{107}

\begin{chapteropening}
\textbf{By the end of this chapter:} \emph{I can say jam, a nut, pepper, a nurse and a cook.}

Complete the last lesson of chapter 107: I can say jam, a nut, pepper, a nurse and a cook.
\end{chapteropening}

\section[\texorpdfstring{\ru{варенье} (varénye)}{варенье (varénye)}]{\ru{варенье} (varénye) — jam}

\label{lesson:RU-C107-jam}



\begin{quote}
Before the new one: say the Russian for sausage, then the Russian for a peach.
\end{quote}

\subsection*{You'll want to know: \ru{варенье}}
\textbf{\ru{варенье}} (\emph{varénye}) — \textquotedblleft{}jam\textquotedblright{}.

\begin{cousinweb}
\textbf{\ru{варенье}} is a noun from \emph{\ru{варить}}, \textquotedblleft{}to boil, to cook\textquotedblright{}: fruit boiled with sugar. Vasmer relates \emph{\ru{варить}} to Lithuanian \emph{virti}, \textquotedblleft{}to boil\textquotedblright{}.
\end{cousinweb}

\begin{grammarlens}[title={What you've built}]
One more word for this chapter.
\end{grammarlens}

\subsection*{Your turn}
\begin{itemize}
\raggedright
  \item \emph{Say it:} \emph{varénye}
  \item \emph{Say it:} \emph{varénye}, once more
  \item \emph{Say it:} say \emph{varénye}
  \item \emph{Recall:} say the Russian for sausage, then the Russian for a peach, then say \emph{varénye} again
\end{itemize}

\begin{tcolorbox}[breakable,colback=teal!4,colframe=teal!35!black,title={Before you move on}]
What does \ru{варенье} mean? (\textquotedblleft{}Jam\textquotedblright{}.) Say it once more.
\end{tcolorbox}
\section[\texorpdfstring{\ru{орех} (orékh)}{орех (orékh)}]{\ru{орех} (orékh) — a nut}

\label{lesson:RU-C107-nut}



\begin{quote}
Before the new one: say the Russian for a peach, then the Russian for jam.
\end{quote}

\subsection*{You'll want to know: \ru{орех}}
\textbf{\ru{орех}} (\emph{orékh}) — \textquotedblleft{}a nut\textquotedblright{}.

\begin{cousinweb}
\textbf{\ru{орех}} $\leftarrow$ Proto-Slavic \emph{*orěxŭ} — inherited. Vasmer compares Lithuanian \emph{ríešutas}, \textquotedblleft{}nut\textquotedblright{}, and Albanian \emph{arrë}, \textquotedblleft{}walnut\textquotedblright{}; Machek saw in the Baltic and Slavic words a loan from an older, unknown language.
\end{cousinweb}

\begin{grammarlens}[title={What you've built}]
One more word for this chapter.
\end{grammarlens}

\subsection*{Your turn}
\begin{itemize}
\raggedright
  \item \emph{Say it:} \emph{orékh}
  \item \emph{Say it:} \emph{orékh}, once more
  \item \emph{Say it:} say \emph{orékh}
  \item \emph{Recall:} say the Russian for a peach, then the Russian for jam, then say \emph{orékh} again
\end{itemize}

\begin{tcolorbox}[breakable,colback=teal!4,colframe=teal!35!black,title={Before you move on}]
What does \ru{орех} mean? (\textquotedblleft{}A nut\textquotedblright{}.) Say it once more.
\end{tcolorbox}
\section[\texorpdfstring{\ru{перец} (pérets)}{перец (pérets)}]{\ru{перец} (pérets) — pepper}

\label{lesson:RU-C107-pepper}



\begin{quote}
Before the new one: say the Russian for jam, then the Russian for a nut.
\end{quote}

\subsection*{You'll want to know: \ru{перец}}
\textbf{\ru{перец}} (\emph{pérets}) — \textquotedblleft{}pepper\textquotedblright{}.

\begin{cousinweb}
\textbf{\ru{перец}} $\leftarrow$ Old Russian \emph{\ru{пьпьрь}}, from 1193 — an old loan from Latin \emph{piper}, which goes back through Greek \emph{péperi} to Sanskrit \emph{pippalī}, \textquotedblleft{}long pepper\textquotedblright{}; Vasmer allows an East Germanic go-between.
\end{cousinweb}

\begin{grammarlens}[title={What you've built}]
One more word for this chapter.
\end{grammarlens}

\subsection*{Your turn}
\begin{itemize}
\raggedright
  \item \emph{Say it:} \emph{pérets}
  \item \emph{Say it:} \emph{pérets}, once more
  \item \emph{Say it:} say \emph{pérets}
  \item \emph{Recall:} say the Russian for jam, then the Russian for a nut, then say \emph{pérets} again
\end{itemize}

\begin{tcolorbox}[breakable,colback=teal!4,colframe=teal!35!black,title={Before you move on}]
What does \ru{перец} mean? (\textquotedblleft{}Pepper\textquotedblright{}.) Say it once more.
\end{tcolorbox}
\section[\texorpdfstring{\ru{медсестра} (medsestrá)}{медсестра (medsestrá)}]{\ru{медсестра} (medsestrá) — a nurse}

\label{lesson:RU-C107-nurse}



\begin{quote}
Before the new one: say the Russian for a nut, then the Russian for pepper.
\end{quote}

\subsection*{You'll want to know: \ru{медсестра}}
\textbf{\ru{медсестра}} (\emph{medsestrá}) — \textquotedblleft{}a nurse\textquotedblright{}.

\begin{cousinweb}
\textbf{\ru{медсестра}} is cut short from \emph{\ru{медицинская} \ru{сестра}}, \textquotedblleft{}medical sister\textquotedblright{}: \emph{\ru{мед}-}, from \emph{\ru{медицина}}, \textquotedblleft{}medicine\textquotedblright{}, with \emph{\ru{сестра}}, \textquotedblleft{}sister\textquotedblright{}, a cousin of English \emph{sister}.
\end{cousinweb}

\begin{grammarlens}[title={What you've built}]
One more word for this chapter.
\end{grammarlens}

\subsection*{Your turn}
\begin{itemize}
\raggedright
  \item \emph{Say it:} \emph{medsestrá}
  \item \emph{Say it:} \emph{medsestrá}, once more
  \item \emph{Say it:} say \emph{medsestrá}
  \item \emph{Recall:} say the Russian for a nut, then the Russian for pepper, then say \emph{medsestrá} again
\end{itemize}

\begin{tcolorbox}[breakable,colback=teal!4,colframe=teal!35!black,title={Before you move on}]
What does \ru{медсестра} mean? (\textquotedblleft{}A nurse\textquotedblright{}.) Say it once more.
\end{tcolorbox}
\section[\texorpdfstring{\ru{повар} (póvar)}{повар (póvar)}]{\ru{повар} (póvar) — a cook}

\label{lesson:RU-C107-cook}



\begin{quote}
Before the new one: say the Russian for pepper, then the Russian for a nurse.
\end{quote}

\subsection*{You'll want to know: \ru{повар}}
\textbf{\ru{повар}} (\emph{póvar}) — \textquotedblleft{}a cook\textquotedblright{}.

\begin{cousinweb}
\textbf{\ru{повар}} is \emph{\ru{по}-} with \emph{\ru{вар}}, from \emph{\ru{варить}}, \textquotedblleft{}to boil, to cook\textquotedblright{}: the one who cooks.
\end{cousinweb}

\begin{grammarlens}[title={What you've built}]
One more word for this chapter.
\end{grammarlens}

\subsection*{Your turn}
\begin{itemize}
\raggedright
  \item \emph{Say it:} \emph{póvar}
  \item \emph{Say it:} \emph{póvar}, once more
  \item \emph{Say it:} say \emph{póvar}
  \item \emph{Recall:} say the Russian for pepper, then the Russian for a nurse, then say \emph{póvar} again
\end{itemize}

\begin{tcolorbox}[breakable,colback=teal!4,colframe=teal!35!black,title={Before you move on}]
What does \ru{повар} mean? (\textquotedblleft{}A cook\textquotedblright{}.) Say it once more.
\end{tcolorbox}
